\documentclass[10pt,a4paper]{amsart}
\usepackage[margin=19mm]{geometry}
\usepackage{amsmath,amssymb,mathtools}
\usepackage[scr=rsfs,cal=euler]{mathalfa}
\usepackage{xfrac}
\usepackage[unicode,hidelinks,bookmarks=false]{hyperref}
\newcommand{\ie}{i.e.\ }
\newcommand{\cf}{cf.\ }
\newcommand{\resp}{respectively\ }
\newcommand{\Subsection}{\subsection}
\providecommand{\nobreakdash}{\nobreak\hbox{-}}
\setlength{\emergencystretch}{2em}
\multlinegap0pt
\usepackage{bm}
\usepackage{bbm}
\usepackage{dsfont}
\usepackage{xspace}
\usepackage{cancel}
\usepackage{mathtools}
\usepackage{amsthm}
\makeatletter
\@ifundefined{theorem}{\newtheorem{theorem}{Theorem}}{}
\@ifundefined{lemma}{\newtheorem{lemma}[theorem]{Lemma}}{}
\@ifundefined{proposition}{\newtheorem{proposition}[theorem]{Proposition}}{}
\makeatother
\DeclareMathOperator{\dom}{dom}
\newcommand\lex{\mathrm{lex}}
\newcommand{\ck}[1]{\check{#1}}
\newcommand{\dt}[1]{\mathring{#1}}
\newcommand{\seb}{\subseteq}
\newcommand{\tle}{\sqsubset}
\renewcommand{\P}{\mathbb{P}}
\title[The class of Aronszajn lines under epimorphisms]{The class of Aronszajn lines under epimorphisms}
\author{Lucas Polymeris, Carlos Martinez-Ranero}
\date{Source: arXiv:2503.13728v2 (CC BY 4.0)}
\begin{document}
\maketitle

\noindent\emph{Source and licence.} The class of Aronszajn lines under epimorphisms. Version arXiv:2503.13728v2 (CC BY 4.0), distributed under the Creative Commons Attribution 4.0 International licence, as recorded in the arXiv licence metadata for this exact version and shown on the abstract page https://arxiv.org/abs/2503.13728v2. Full rights notice: licenses/arxiv-2503.13728-CC-BY-4.0.txt.\par\smallskip

\noindent\textit{Editorial scope.} Counted unit: the Proposition whose source label is \texttt{None} in the article (source0.tex lines 2343--2345, source label \texttt{None}), delivered together with the article's own complete original proof of it (source0.tex lines 2346--2409, one \texttt{proof} environment of 64 lines). No mathematics is written, repaired or re-derived: statement and proof are the article's own text line for line. Required context retained uncounted: lem:halbeisen (lines 2328--2333). Every definition, display and label of those lines is retained unchanged. Omitted from this package and not redistributed: 1--2327, 2334--2342, 2410--2608. Nothing inside a retained excerpt is omitted or altered: every retained body is delivered exactly as the article's lines are written, so no omission receipt of any kind accompanies this package. Citations: the retained excerpts carry no citation command at all, so no bibliography entry of the article is transcribed and no reference is redistributed. Reference adaptation: none was needed and none was made; every reference inside the delivered unit resolves inside it, so no unresolved reference or citation is produced. Printed slips kept literal: no printed word, symbol, hypothesis or number of the counted statement or of its proof was altered. Layout and class: the article's own class is \texttt{amsart} with packages including lmodern, fontenc, inputenc, amsmath, amssymb, mathrsfs, comment, color, cite, hyperref, enumitem, eucal, tikz, tikzscale; the wrapper is the lane's pinned \texttt{amsart} editorial wrapper at 10pt on A4, which restates the theorem-like environments the retained text needs and adds the article's own macro definitions that the retained text needs (\texttt{\textbackslash P}, \texttt{\textbackslash ck}, \texttt{\textbackslash dom}, \texttt{\textbackslash dt}, \texttt{\textbackslash lex}, \texttt{\textbackslash seb}, \texttt{\textbackslash tle}), with extra wrapper packages (bm, bbm, dsfont, xspace, cancel, mathtools). The delivered numbering is the unit's own. Assets: no retained span contains a figure environment, a graphics-inclusion command, a PDF-page inclusion, a table or a TikZ picture; no image file is copied into this package, the package declares no asset dependency and no image file is redistributed with it. Licence: version arXiv:2503.13728v2 of this article is distributed under the Creative Commons Attribution 4.0 International licence, as recorded in the arXiv licence metadata for this exact version and shown on the abstract page https://arxiv.org/abs/2503.13728v2. A full rights notice is delivered at licenses/arxiv-2503.13728-CC-BY-4.0.txt. Characters: every retained excerpt is pure ASCII, so the delivered engine, XeLaTeX with its default Latin Modern text font, reports no missing character.
\begin{lemma}\label{lem:halbeisen}
  Let $K \seb T$ be uncountable, and let $n \in \omega$. Then
  there is an uncountable $K' \seb K$ such that
  for all $s,t \in K'$, if $\xi \in \dom(s) \cap \dom(t)$,
  and $t(\xi) < n$ or $s(\xi) <n$, then $t(\xi) = s(\xi)$.
\end{lemma}

\begin{proposition}
  $V[G] \models$ $(T^{c},<_{\lex})$ is $\omega_{1}$-irreversible.
\end{proposition}

\begin{proof}
  Suppose towards a contradiction
  that $V[G] \models T^{c} \text{ is not $\omega_{1}$-irreversible}$.
  Then there must be $p \in G$ and names $\dt{A}$,$\dt{f}$ such
  that
  \[p \Vdash \dt{A} \text{ is uncountable and }
    f : \dt{A} \seb T^{\dt{c}} \to T^{\dt{c}} \text{ reverses } <_{\lex}.\]
  To arrive to a contradiction it is enough to find some extension
  of $p$ that forces
  $(\exists t,s \in \dt{A})(t <_{\lex} s \land \dt{f}(t) <_{\lex} \dt{f}(s))$.

  For each $q \le p$, consider the
  set $A_{q} := \{t \in T : q \Vdash \dt{c} \circ \ck{t} \in \dt{A}\}$. Since $\P$ is
  countable, there is some $q \le p$ such that $A_{q}$ is uncountable. Fix such a $q$.
  In similar fashion, there is $q' \le q$, an uncountable $K \seb A_{q}$,
  and $g : K \to T$ in $V$ such that for all $t \in K$,
  $q' \Vdash \dt{f}(\dt{c} \circ \ck{t}) = \dt{c} \circ \ck{g}(\ck{t})$. We may
  assume that this is already forced by $q$. Let $n$ be greater
  than any element in $\dom(q)$. By going to an uncountable
  subset of $K$, we may assume that $K$ and $\{g(t) : t \in K\}$ satisfy
  the conclusions of Lemma~\ref{lem:halbeisen}.

  Since $T$ is Countryman we see that $T$ is
  $\omega_{1}$-irreversible, and thus there are $t,s \in K$
  such that $t <_{\lex} s$ and $g(t) <_{\lex} g(s)$. There are four
  cases depending on the extension relation of $t$ with $s$ and
  of $g(t)$ with $g(s)$. In
  case that $t \not\tle s$ we let $x$ and $y$ be such that
  $t(\xi) = x$ and $s(\xi) = y$, where $\xi$ is the first such that
  $t(\xi) \neq s(\xi)$. Otherwise we let $x$ and $y$ undefined.
  In similar fashion, when $g(t) \not\tle g(s)$
  we let $a = g(t)(\xi)$ and $b = g(s)(\xi)$ where
  $\xi$ is the first such that $g(t)(\xi) \neq g(s)(\xi)$.
  Note that when they are defined,
  $a,b,x,y > n$, $x < y$ and $a < b$.

  If $t \tle s$ and $g(t) \tle g(s)$, then clearly
  $q \Vdash \dt{c} \circ \ck{t} \tle \dt{c} \circ \ck{s}$ and
  $q \Vdash \dt{c} \circ \ck{g}(\ck{t}) \tle \dt{c} \circ \ck{g}(\ck{s})$.
  This arrives gives us our desired contradiction, since
  $q \Vdash \dt{c} \circ \ck{t}, \dt{c} \circ \ck{s} \in \dt{A}$ and
  $q \Vdash \dt{c} \circ \ck{g}(\ck{t}) = \dt{f}(\dt{c} \circ \ck{t}) \land
  \dt{c} \circ \ck{g}(\ck{s}) = \dt{f}(\dt{c} \circ \ck{s})$.

  If $t \tle s$ and $g(t) \not\tle g(s)$, then
  letting $q' := q \cup \{(a,0),(b,1)\}$ we see that $q' \in \P$,
  $q' \le q$ and
  $q' \Vdash \dt{c} \circ \ck{t} \tle \dt{c} \circ \ck{s} \land
  \dt{c} \circ \ck{g}(\ck{t}) <_{\lex} \dt{c} \circ \ck{g}(\ck{s})$,
  which is again a contradiction.
  If $t \not\tle s$ and $g(t) \tle g(s)$ the argument is identical.

  Assume now that $t \not\tle s$ and $g(t) \not\tle g(s)$.
  If $x = b$,
  then $a < x = b < y$ and thus
  $q' := q \cup \{(a,0),(x,1),(y,2)\}$ is in $\P$. But note then
  $q' \le q$ and
  $q \Vdash \dt{c} \circ \ck{t} <_{\lex} \dt{c} \circ \ck{s}
  \land \dt{c} \circ \ck{g}(\ck{t}) <_{\lex} \dt{c} \circ \ck{g}(\ck{s})$,
  and we arrive at a contradiction as before. Assume now $x \neq b$.
  If $y = a$
  then $q' := q \cup \{(x,0),(a,1),(b,2)\}$ works, and if $y \neq a$
  then $q' := q \cup \{(x,0),(y,1),(a,0),(b,1)\}$ works.
\end{proof}

\end{document}
